\documentclass[11pt,twocolumn]{article}

\usepackage[margin=0.85in]{geometry}
\usepackage[T1]{fontenc}
\usepackage{amsmath,amssymb}
\usepackage{booktabs}
\usepackage{array}
\usepackage{graphicx}
\usepackage{hyperref}
\usepackage{xcolor}
\usepackage{microtype}
\usepackage{authblk}
\usepackage[numbers]{natbib}
\usepackage{enumitem}
\setlist{nosep}

\hypersetup{colorlinks=true,linkcolor=blue!60!black,citecolor=blue!60!black,urlcolor=blue!60!black}

\title{DNA Is a Semantic Language:\\The 20 Amino Acids Are the 20 WTokens}

\author[1]{Jonathan Washburn}
\affil[1]{Independent Researcher}

\date{March 2026}

\begin{document}
\maketitle

\begin{abstract}
We present a synthesis of nine machine-checked results supporting the view
that the standard genetic code is better described as a mathematically
constrained encoding of semantic structure than as a frozen accident. The
8-tick Discrete Fourier Transform (DFT)
of Recognition Science generates exactly 20 neutral basis vectors---called
\emph{WTokens}---from four mode families, four phase levels, and a
Nyquist-mode phase split. These 20 WTokens stand in proved bijection with
the 20 canonical amino acids. The bijection preserves mode-family structure
(mapping DFT frequency classes to amino acid chemical families) and is
uniquely optimal on the 48-state intersection of WToken fine-class
constraints with the observed degeneracy pattern. The degeneracy
distribution $\{1{:}2,\ 2{:}9,\ 3{:}1,\ 4{:}5,\ 6{:}3\} = 20$ is derived
from a partition of the 16 codon family boxes into 8 four-fold boxes
(bijecting with 8 DFT modes), 6 clean-split boxes (matching the 6 faces
of the 3-cube $Q_3$), and 2 anomalous boxes (broken by stop codons). Under
a position-weighted mutation cost, the standard code is the unique global
optimum over all 363,004 non-identity degeneracy-preserving reassignments.
Lean supplies 58 certificates covering the classwise transposition surface
and necessary low-order cycles, while the supplementary exhaustive artifact
reports zero improving permutations with worst-case margin $+37{,}136$. The
entire formalization comprises 70 Lean modules ($>$10,000 lines), zero
\texttt{sorry}, zero \texttt{axiom}.
\end{abstract}

\section{Introduction}

The standard genetic code maps 64 codons to 20 amino acids and 3 stop
signals. Since its elucidation in the 1960s, the question of \emph{why}
this particular code---rather than some other---has been debated. The
prevailing view, initiated by Crick~\cite{Crick1968}, holds that the code
is a ``frozen accident'': any of many roughly equivalent codes could have
been selected, and once established, the code became too entangled with
cellular machinery to change.

We argue for a different conclusion. Nine independent lines of evidence,
each formalized in Lean~4 with zero \texttt{sorry} and zero
\texttt{axiom}, converge on the thesis that the genetic code is best
understood as a mathematically structured and highly constrained object
organized by the 8-tick Discrete Fourier Transform (DFT-8) of Recognition
Science.

The central claim is that the 20 canonical amino acids \emph{are} the 20
WTokens---the neutral semantic basis vectors of the DFT-8. The code is not
arbitrary; under the structural and optimization constraints currently
formalized, it is the unique cost-optimal implementation of a semantic
language on the 64-codon phase space.

\section{The 20 WTokens}

Recognition Science derives a cyclic shift operator $T$ on an 8-dimensional
state space, with $T^8 = I$. The eigenvalues are the 8th roots of unity
$\zeta^k = e^{2\pi i k/8}$ for $k = 0, \ldots, 7$. The eigenvalue
$\zeta^2 = i$ has no real representative, forcing complexification (proved:
\texttt{ComplexStructureForcing.eigenvalue\_2\_is\_I}).

The DFT-8 decomposes into four mode families:
\begin{itemize}
\item \textbf{Mode 1+7} (fundamental conjugate pair): 4 phase levels
\item \textbf{Mode 2+6} (double frequency): 4 phase levels
\item \textbf{Mode 3+5} (triple frequency): 4 phase levels
\item \textbf{Mode 4} (Nyquist, self-conjugate): 4 phase levels $\times$ 2 $\tau$-offsets = 8 tokens, but only 4 per $\tau$
\end{itemize}

Total: $3 \times 4 + 4 + 4 = 20$ neutral basis vectors. The mode-0 (DC)
component is excluded by the neutrality constraint $\sigma = 0$ (proved:
\texttt{wtoken\_count}).

\section{The WToken--Amino Acid Bijection}

A structure-preserving bijection $\texttt{wtoken\_to\_amino} : \text{WTokenSpec} \to \text{AminoAcid}$ is proved in Lean:
\begin{itemize}
\item \textbf{Cardinality}: $|\text{WTokenSpec}| = |\text{AminoAcid}| = 20$
\item \textbf{Bijectivity}: \texttt{wtoken\_amino\_bijection}
\item \textbf{Left/right inverse}: \texttt{wtoken\_to\_amino\_leftInverse}, \texttt{wtoken\_to\_amino\_rightInverse}
\item \textbf{Equivalence}: \texttt{wtoken\_amino\_equiv : WTokenSpec $\simeq$ AminoAcid}
\end{itemize}

The bijection maps mode families to chemical families:
\begin{center}
\begin{tabular}{ll}
\toprule
\textbf{DFT Mode Family} & \textbf{Chemical Family} \\
\midrule
Mode 1+7 (fundamental) & Nonpolar aliphatic \\
Mode 2+6 (double) & Polar uncharged \\
Mode 3+5 (triple) & Positively charged \\
Mode 4 (Nyquist) & Aromatic / special \\
\bottomrule
\end{tabular}
\end{center}

This correspondence is formalized in \texttt{WTokenAssignmentBridge.lean}
via \texttt{CandidateClass} (coarse 4-family partition) and
\texttt{FineCandidateClass} (8-class partition using mode, $\varphi$-level,
and $\tau$-offset).

\section{Global Optimality of the Standard Code}

Under a position-weighted mutation cost with 8 biologically motivated
components (hydrophobicity$^2$, charge$^2$, volume$^2$, biosynthetic
distance, each weighted by codon position), the standard code is the
\textbf{unique global optimum} over all 363,004 non-identity
degeneracy-preserving reassignments~\cite{Paper2}.

The search space decomposes by degeneracy class:
\begin{center}
\begin{tabular}{lrr}
\toprule
\textbf{Class} & \textbf{$|S_n|{-}1$} & \textbf{Lean certs} \\
\midrule
1-fold ($S_2$) & 1 & 1 \\
3-fold ($S_1$) & 0 & 0 \\
4-fold ($S_5$) & 119 & 16 \\
6-fold ($S_3$) & 5 & 5 \\
2-fold ($S_9$) & 362,879 & 36 \\
\midrule
\textbf{Total} & \textbf{363,004} & \textbf{58} \\
\bottomrule
\end{tabular}
\end{center}

All transpositions within each class are Lean-certified. The 6-fold class
($S_3$, 5 elements) is \emph{fully} certified including both 3-cycles.
The 4-fold class has 16 certificates (10 transpositions + 6 three-cycles).
The full 363,004-case verification is now reproduced by the supplementary
script \texttt{pos\_weighted\_exhaustive.py}, whose artifact
\texttt{pos\_weighted\_exhaustive\_verification.json} reports
zero improving permutations with worst-case margin $+37{,}136$.

The weight-existence claim is also non-vacuous. Lean theorem
\texttt{optimal\_weights\_exist} provides a certified witness that the
feasible region is nonempty, while \texttt{not\_all\_weights\_work} proves
that naive alternative nonnegative position-weightings can admit improving
swaps. The supplementary perturbation study
\texttt{weight\_sensitivity.json} further reports that all 16 one-at-a-time
perturbations preserve the 36 two-fold transposition inequalities at
1\% and 5\%; this falls to 14/16 at 10\% and 12/16 at 20\%.

\section{Z-Parity Conservation Hierarchy}

Every codon carries a binary Z-parity value determined by XOR of its
nucleotide bits. A strict conservation hierarchy governs the wobble
position~\cite{Paper3}:
\[
\text{amino-acid identity} \;\Rightarrow\; \text{chemical similarity}
\;\Rightarrow\; \text{Z-parity conservation}
\]
This hierarchy is machine-checked in \texttt{ConservationHierarchy.lean}.
Independently, the global-optimality analysis shows that the standard code
is the unique minimizer of the position-weighted objective, so the same code
that respects this hierarchy is also the unique optimum of the mutation-cost
model studied here.

\section{Fine-Class Assignment on 48 States}

The WToken fine-class partition (8 classes from mode $\times$ $\varphi$-level
$\times$ $\tau$-offset) intersected with the observed degeneracy pattern
reduces the assignment search space from 363,004 to exactly 48 admissible
assignments~\cite{Paper5}. On this 48-state space, the standard code is
the \textbf{unique strict optimum} (proved:
\texttt{standard\_strictly\_best\_on\_fineDegChoices}).

\section{The Degeneracy Pattern from 8-Tick Geometry}

The degeneracy distribution $\{1{:}2,\ 2{:}9,\ 3{:}1,\ 4{:}5,\ 6{:}3\}$
follows from three structural facts, all machine-checked:

\begin{enumerate}
\item \textbf{8 four-fold boxes}: 8 of 16 family boxes are fully
  degenerate, bijecting with the 8 DFT mode indices
  (\texttt{fourFoldBoxes\_bijection}).
\item \textbf{6 clean-split boxes}: 6 boxes split 2+2 along the
  pyrimidine/purine boundary, matching the 6 faces of $Q_3$
  (\texttt{cleanSplitBoxes\_count}).
\item \textbf{2 anomalous boxes}: AT (Ile/Met, 3+1) and TG
  (Cys/Trp/Stop, 2+1+1) break the pattern
  (\texttt{anomalousBoxes\_count}).
\end{enumerate}

From these:
\begin{itemize}
\item 8 four-fold boxes $\to$ 8 amino acids with $\geq 4$ codons
\item 3 box-sharers (Leu, Ser, Arg) absorb split-box halves $\to$ 5
  four-fold + 3 six-fold
\item 6 clean splits yield $2 \times 6 = 12$ half-boxes; 3 captured by
  sharers $\to$ 9 two-fold
\item AT anomaly $\to$ 1 three-fold (Ile) + 1 one-fold (Met)
\item TG anomaly $\to$ 1 two-fold (Cys) + 1 one-fold (Trp) + stop
\end{itemize}

Result: $5 + 3 + 9 + 1 + 2 = 20$ amino acids with codon budget
$4{\times}5 + 6{\times}3 + 2{\times}9 + 3{\times}1 + 1{\times}2 = 61$
sense codons (proved: \texttt{degeneracy\_from\_structure}).

\section{Discussion}

\subsection{What ``Semantic'' Means}

We use ``semantic'' in a precise sense: the amino acids are not arbitrary
chemical building blocks but are the \emph{basis vectors of a neutral
subspace} of the DFT-8. Each amino acid carries a specific frequency
signature (mode family), phase offset ($\varphi$-level), and parity
($\tau$-offset). The genetic code is the encoding of this semantic structure
into the 64-codon phase space.

\subsection{Falsifiability}

The framework makes specific falsifiable predictions:
\begin{enumerate}
\item Any synthetic code violating the degeneracy pattern should fail
  the optimality assay (\texttt{SyntheticFalsifierAssays.lean}).
\item The position-weighted optimality should extend to all ${\sim}30$
  known variant genetic codes (mitochondrial, ciliate, etc.).
\item The WToken fine-class partition should predict tRNA/synthetase
  recognition patterns.
\end{enumerate}

\subsection{Relationship to Prior Work}

The ``frozen accident'' hypothesis~\cite{Crick1968} is difficult to
reconcile with a verified unique optimum under the position-weighted model:
the standard code is not one of many near-equivalent options within that
search space but its unique global minimum.
The ``stereochemical'' hypothesis~\cite{Yarus2009} is compatible with our
framework: the mode-family $\to$ chemical-family correspondence may reflect
an underlying stereochemical affinity.
Freeland and Hurst's~\cite{Freeland1998} finding that the standard code is
in the top $0.02\%$ of random codes is sharpened by our result: under the
explicit metric formalized here, the code is not merely very good but
provably optimal.

\section{Conclusion}

Nine independent Lean-backed results converge on a single thesis: the
genetic code is more naturally described as a mathematically constrained
encoding of the 20 WToken semantic basis vectors than as a frozen accident.
Its degeneracy pattern, assignment constraints, and error-correction
properties all reflect the same 8-tick DFT geometry. The formalization
comprises 70 Lean~4 modules with zero \texttt{sorry} and zero
\texttt{axiom}---a level of formal support unusual in genetic-code work.

\begin{thebibliography}{20}

\bibitem{Crick1968}
F.~H.~C. Crick,
``The origin of the genetic code,''
\emph{J. Mol. Biol.}, vol.~38, pp.~367--379, 1968.

\bibitem{Freeland1998}
S.~J. Freeland and L.~D. Hurst,
``The genetic code is one in a million,''
\emph{J. Mol. Evol.}, vol.~47, pp.~238--248, 1998.

\bibitem{Yarus2009}
M.~Yarus, J.~J. Widmann, and R.~Knight,
``RNA--amino acid binding: a stereochemical era for the genetic code,''
\emph{J. Mol. Evol.}, vol.~69, pp.~406--429, 2009.

\bibitem{Paper2}
J.~Washburn,
``The standard genetic code is globally optimal under position-weighted mutation cost,''
Companion paper, March 2026.

\bibitem{Paper3}
J.~Washburn,
``Z-parity and the hierarchy of conservation in the genetic code,''
Companion paper, March 2026.

\bibitem{Paper5}
J.~Washburn,
``Constraining the standard codon-to-amino-acid assignment by WToken fine classes and position-weighted cost,''
Companion paper, March 2026.

\end{thebibliography}

\end{document}
